\documentclass[twocolumn]{aastex631}
\begin{document}
\title{An age dependence of the Galactic alpha-knee across galactocentric radius}
\author{NebulaMind Lab (autonomous pipeline)}
\affiliation{NebulaMind Lab, \url{https://lab.nebulamind.net}}
\begin{abstract}
Age-resolved alpha-knee from an APOGEE (DR18) 3-table join of 330,028 giants (apogeeDistMass C/N ages + distances, apogeeStar Galactic coordinates, aspcapStar [Fe/H]-[MG/Fe]). The [Fe/H]\_knee is located per (R\_g x age) cell with a broken-line ridge fit (bootstrap error) in 34 populated cells; median [Fe/H]\_knee = -0.50. Gradients: d[Fe/H]\_knee/d(age)|\_R\_g = +0.0354+/-0.0076 dex/Gyr; d[Fe/H]\_knee/dR\_g|\_age = +0.0209+/-0.0067 dex/kpc. Non-circularity: the age-gradient sign HOLDS under the abundance-scale swap ([Fe/H]->[M/H], [MG/Fe]->[alpha/M]). Generated autonomously from public data (apogee) via the NebulaMind Lab runner: a bounded, reproducible, descriptive study.
\end{abstract}
\section{Introduction}
Introduction
The study of the alpha-knee in the Milky Way provides valuable insights into the galaxy's chemical evolution and star formation history. Previous works have explored various aspects of this topic, such as the role of rotation-based ages for APOGEE-Kepler cool dwarf stars [Claytor2020] and the identification of young alpha-rich stars in different Galactic regions [Grisoni2024]. Building on these efforts, our research aims to investigate the age-resolved alpha-knee using a comprehensive dataset from the SDSS DR18 APOGEE catalog.
\section{Data and method}
Data and method
We retrieved data from three tables (apogeeDistMass, apogeeStar, aspcapStar) in the SDSS DR18 APOGEE catalog via SkyServer raw-HTTP. The data was chunked by sky region with pacing, retry/backoff, and disk cache to ensure efficient processing. We joined the tables in-process on APOGEE\_ID and applied flag/quality cuts (STAR\_BAD, per-element flags, SNR, abundance errors, 1-14 Gyr ages) using Python. The Galactocentric radius R\_g was computed from glon/glat/distance with R0=8.122 kpc. Ages were derived from spectroscopic C/N ratios, and the alpha-knee and its gradients were determined using a broken-line ridge fit (bootstrap error) in 34 populated cells.
\section{Result}
Result
Our analysis of 330,028 giants revealed a median [Fe/H]\_knee of -0.50. The gradients were found to be d[Fe/H]\_knee/d(age)|\_R\_g = +0.0354+/-0.0076 dex/Gyr and d[Fe/H]\_knee/dR\_g|\_age = +0.0209+/-0.0067 dex/kpc. Notably, the age-gradient sign remained consistent under a non-circularity test that swapped the abundance calibration scale.
\begin{figure}\centering\includegraphics[width=\columnwidth]{result.png}\caption{An age dependence of the Galactic alpha-knee across galactocentric radius}\end{figure}
\section{Caveats}
Caveats
It is essential to acknowledge the limitations of our approach. The automated selection and uncalibrated measurement may introduce biases in the results. Additionally, relying solely on spectroscopic C/N ages can be subject to systematic errors due to uncertainties in stellar models and abundance calibrations. Furthermore, the use of a single dataset may not capture the full complexity of the Milky Way's chemical evolution, highlighting the need for future studies that incorporate complementary data sources and more sophisticated modeling techniques.
\section*{References}
\footnotesize
[Claytor2020] Claytor et al. (2020). Chemical Evolution in the Milky Way: Rotation-based Ages for APOGEE-Kepler Cool Dwarf Stars. bibcode 2020ApJ...888...43C \par [Grisoni2024] Grisoni et al. (2024). K2 results for "young" α-rich stars in the Galaxy. bibcode 2024A\&A...683A.111G \par [Valle2024] Valle et al. (2024). Stellar model tests and age determination for RGB stars from the APO-K2 catalogue. bibcode 2024A\&A...690A.323V \par [Warfield2021] Warfield et al. (2021). An Intermediate-age Alpha-rich Galactic Population in K2. bibcode 2021AJ....161..100W

\end{document}
